% Pregunta 71 (2024-junio-A, n.º 24): regla de oro de Phelps en Solow con f(k)=k^0,3, n+δ=0,15. Autor: Víctor Gutiérrez Marcos
\documentclass[tikz,border=6pt]{standalone}
\input{tikz-tcee}
\begin{document}
\begin{tikzpicture}[tcee,yscale=2.5]
  % ejes (en y se escala 2,5: 1 unidad de y = 2,5 cm)
  \draw[eje] (0,0) -- (7.2,0) node[below left=2pt and 0pt] {$k$};
  \draw[eje] (0,0) -- (0,2.05) node[left] {$y$};
  \node[below left] at (0,0) {$0$};
  % s = 0,5: estado estacionario k = (0,5/0,15)^{1/0,7} ~ 5,58
  \draw[curvanueva,desplazada,domain=0.002:6.6,samples=120] plot (\x,{0.5*\x^0.3}) node[etiqueta,below right=0pt and 0pt,tceerojo] {$0{,}5\,f(k)$};
  % f(k)
  \draw[curva1,domain=0.002:6.6,samples=120] plot (\x,{\x^0.3}) node[etiqueta,right,tceemorado] {$f(k)=k^{0{,}3}$};
  % (n+δ)k
  \draw[curva2,domain=0:6.6] plot (\x,{0.15*\x}) node[etiqueta,above right=0pt and 0pt,tceeazul] {$(n+\delta)k$};
  % s_oro f(k)
  \draw[curva3,domain=0.002:6.6,samples=120] plot (\x,{0.3*\x^0.3}) node[etiqueta,right,tceeverde] {$0{,}3\,f(k)$};
  % tangente en k_oro, pendiente 0,15
  \draw[tceedorado,thick,dashed] (0.9,{1.3459+0.15*(0.9-2.6918)}) -- (4.6,{1.3459+0.15*(4.6-2.6918)}) node[etiqueta,above,tceedorado,align=center] {pendiente\\$f'(k)=n+\delta$};
  % c_oro
  \draw[guia] (2.6918,0) -- (2.6918,0.4038);
  \draw[<->,very thick,tceemorado] (2.6918,0.4038) -- (2.6918,1.3459) node[midway,left,etiqueta,tceemorado] {$c_{\mathrm{oro}}$};
  \node[punto] at (2.6918,0.4038) {};
  \node[punto] at (2.6918,1.3459) {};
  \node[below] at (2.6918,0) {$k_{\mathrm{oro}}\approx2{,}69$};
  % c con s=0,5
  \draw[guia] (5.583,0) -- (5.583,0.8374);
  \draw[<->,thick,tceegris] (5.583,0.8374) -- (5.583,1.6748) node[midway,left,etiqueta,tceegris] {$c$};
  \node[punto] at (5.583,0.8374) {};
  \node[below] at (5.583,0) {$k^*\approx5{,}58$};
\end{tikzpicture}
\end{document}
